% ===================================================================
%  TABULKA č. 3 — Dětství a mládí.
%  Traduction 2026 d'apres texto/io, controlee sur texto/fr, texto/ru
%  et texto/uk. Consigne : voir texto/cs/10-tabulka-01.tex.
%
%  LES MOIS ET LES JOURS SEPARENT LE TCHEQUE DE TOUTES LES AUTRES
%  COLONNES DU VOLUME, y compris de sa voisine slave, et c'est le
%  seul cas de ce genre en treize colonnes. Vingt-cinq langues disent
%  le latin — mars, avril, mai, ou leurs formes — et le russe comme
%  l'ukrainien ont leurs propres noms slaves. Le tcheque en a
%  d'autres encore : « březen » le mois des bouleaux, « duben » le
%  mois des chenes, « květen » le mois des fleurs. Ni « марта » ni
%  « березень » ne s'y lisent.
%
%  LES JOURS FONT DE MEME, ET ILS COMPTENT. « Pondělí » c'est
%  l'apres-dimanche, « úterý » le deuxieme, « středa » le milieu,
%  « čtvrtek » le quatrieme, « pátek » le cinquieme. Le portugais
%  comptait aussi, au tableau 3 galicien, mais a partir du dimanche
%  et par un autre mot. Deux langues qui comptent, et pas les memes
%  nombres.
%
%  L'ENTREE DE JUILLET EST DECLAREE DANS gravuri/korekti.json, comme
%  pour les vingt-six colonnes precedentes : le fac-simile ido
%  compose « (julio od agosto) » avec « agosto » en gras et « julio »
%  en romain, alors que les deux sont des mois. LA SOURCE NE BOUGE
%  PAS ; c'est la lecture qui se corrige, langue par langue.
% ===================================================================

%%K t03-apar-1 apar
\VUsaut{3.0mm}
\VUtitre{15.0pt}{60}{TABULKA č. 3}
\VUsaut{1.6mm}
\VUtitre{11.4pt}{0}{Dětství a mládí.}
\VUsaut{3.4mm}

%%K t03-apar-2 apar
(Tabulky 3 a 4 jsou sestaveny tak, že ve čtyřech \VUgras{rodinných výjevech} předkládají
základní pojmy o \VUgras{počasí}, \VUgras{věku}, \VUgras{rodině},
\VUgras{oblečení} a \VUgras{stavu}.)
\VUsaut{4.0mm}

%%K t03-tit-1 sub
\VUcentre{12.0pt}{0}{\textit{První scéna.}}
\VUsaut{3.0mm}

%%K t03-tit-2 sub
\VUpk{10.2pt}{40}{Křtiny na vesnici.}
\VUsaut{2.4mm}

%%K t03-tit-3 sub
\VUpk{10.2pt}{40}{Jaro.}
\VUsaut{3.0mm}

%%K t03-c1-01-1 p
1. --- Jsme na \VUgras{jaře}, v měsících \VUgras{březnu}, \VUgras{dubnu} nebo
\VUgras{květnu}, ve dnech \VUgras{týdne}, v \VUgras{pondělí}, \VUgras{úterý} nebo
\VUgras{středu}. Je \VUgras{deset hodin} ráno.

%%K t03-c1-02-1 p
2. --- \VUgras{Počasí} je krásné, \VUgras{vzduch} čistý. Slunce svítí. Několik
\VUgras{obláčků} \textsuperscript{(2)} stoupá na modrém \VUgras{nebi}
\textsuperscript{(1)}.

%%K t03-c1-03-1 p
3. --- \VUgras{Stromy} mají sotva \VUgras{listy}. Štíhlé \VUgras{topoly} \textsuperscript{(3)}
a vysoký \VUgras{platan} \textsuperscript{(17)} mají zatím jen pupeny.

%%K t03-c1-04-1 p
4. --- \VUgras{Vlaštovky} \textsuperscript{(4)} se vracejí a poletují s radostným štěbetáním
kolem \VUgras{špičky} \textsuperscript{(7)} \VUgras{zvonice} \textsuperscript{(5)}, která
korunuje vesnický \VUgras{kostel} \textsuperscript{(6)}.

%%K t03-tit-4 sub
\VUsaut{3.4mm}
\VUpk{10.2pt}{40}{Kostel.}
\VUsaut{3.0mm}

%%K t03-c1-05-1 p
5. --- Velmi prostý je ten kostel se svým velkým \VUgras{portálem} \textsuperscript{(13)} a
mohutnými \VUgras{hodinami} \textsuperscript{(8)}. Malá
\VUgras{ručička} \textsuperscript{(9)} je na číslici X; ukazuje \VUgras{hodiny.}
\VUgras{Velká} \textsuperscript{(10)} je na XII (poledne); ukazuje \VUgras{minuty.}

%%K t03-c1-06-1 p
6. --- Ve zvonici vesele znějí \VUgras{zvony} \textsuperscript{(11)}. Na vrcholu střechy stojí
malý kamenný \VUgras{kříž} \textsuperscript{(12)}.

%%K t03-c1-07-1 p
7. --- Kostel nemá jen velký \VUgras{portál} \textsuperscript{(13)}; má také malá boční
dvířka. Těmi dvířky vychází pan \VUgras{farář} \textsuperscript{(15)} oblečený ve své
\VUgras{sutaně} ze \VUgras{sakristie} \textsuperscript{(14)} a jde na
\VUgras{faru} \textsuperscript{(19)}.

%%K t03-c1-08-1 p
8. --- Blízko něho, zde je \VUgras{mlékařka} \textsuperscript{(24)} se svým
\VUgras{oslem} \textsuperscript{(23)}. Vrací se z města, kam přinesla dobré mléko pro
děti.

%%K t03-tit-5 sub
\VUsaut{4.0mm}
\VUpk{10.2pt}{40}{Ovocný sad.}
\VUsaut{3.4mm}

%%K t03-c1-09-1 p
9. --- Pan Aliboron (osel) je s tou ranní cestou zcela spokojen. S rozkoší vdechuje vzduch
provoněný vůní \VUgras{květů} \textsuperscript{(20)}, které pokrývají
\VUgras{ovocné stromy} \textsuperscript{(18)} sousedního \VUgras{ovocného sadu.}
Zcela růžové a bílé jsou ty \VUgras{broskvoně} \textsuperscript{(18)} a
\VUgras{meruňky} \textsuperscript{(19)} a dávají naději na dobrou sklizeň.

%%K t03-c1-10-1 p
10. --- Zatím malé \VUgras{včely} \textsuperscript{(22)} z \VUgras{úlu} \textsuperscript{(21)}
tam letí na lup sbírat svůj sladký \VUgras{med} a bzučí.

%%K t03-c1-11-1 p
11. --- Ale co se děje? To vesnické děti přibíhají a plaší polekané \VUgras{husy}
\textsuperscript{(25)}. Je to východ průvodu z kostela po \VUgras{křtu} notářova syna.

%%K t03-c1-12-1 p
12. --- Malý Jakub se ve své \VUgras{kolébce} \textsuperscript{(26)} probouzí radostným
zvoněním zvonů, a jeho sestra Magdalena, která také chce vidět křestní průvod, prosí matku,
aby jí přišla učesat vlasy a obléci ji na prahu domu.

%%K t03-c1-13-1 p
13. --- Matka svoluje. Vezme si svůj \VUgras{šátek} \textsuperscript{(28)}, svůj
\VUgras{živůtek} \textsuperscript{(29)} a svou \VUgras{zástěru}
\textsuperscript{(30)} a přichází si sednout na malou židli přede dveře. Magdalena má jen svou
\VUgras{spodničku} \textsuperscript{(31)} a svůj \VUgras{šněrovací živůtek}
\textsuperscript{(32)}.

%%K t03-c1-14-1 p
14. --- Jak je šťastná! Zapomíná na své oblíbené květiny, na své \VUgras{karafiáty}
\textsuperscript{(34)}, na které usedají \VUgras{motýli} \textsuperscript{(35)}, na své
\VUgras{tulipány} \textsuperscript{(36)}, na své \VUgras{konvalinky} \textsuperscript{(37)} v
\VUgras{květináčích} \textsuperscript{(38)} na \VUgras{zdi} \textsuperscript{(33)} a na své
\VUgras{růže} \textsuperscript{(39)}, které tvoří tak krásný plot. Prohlíží si bohaté oblečení
dam z průvodu.

%%K t03-c1-15-1 p
15. --- Vesnické děti si oblečení příliš nevšímají. Ženou se a strkají do sebe, aby dostaly
\VUgras{bonbony} \textsuperscript{(55)} nebo \VUgras{mince} \textsuperscript{(52)}. Jedno z
nich pláče \textsuperscript{(40)}.

%%K t03-c1-16-1 p
16. --- Druhé šláplo na tlapu \VUgras{psu} \textsuperscript{(42)}, který s vytím utíká a plaší
\VUgras{slepici} \textsuperscript{(43)} a její \VUgras{kuřátka} \textsuperscript{(44)}.

%%K t03-tit-6 sub
\VUsaut{5.0mm}
\VUpk{10.2pt}{40}{Křestní průvod.}
\VUsaut{4.0mm}

%%K t03-c1-17-1 p
17. --- Před průvodem kráčí \VUgras{chůva} \textsuperscript{(45)}. Má krásný
\VUgras{čepec} \textsuperscript{(46)} s dlouhými stuhami. Nese \VUgras{novorozeně}
\textsuperscript{(47)} celé oblečené do \VUgras{krajek} \textsuperscript{(48)}.

%%K t03-c1-18-1 p
18. --- Za chůvou jdou \VUgras{kmotr} \textsuperscript{(49)} a \VUgras{kmotra}
\textsuperscript{(53)}. Rozdávají bonbony a mince. Kmotr má \VUgras{rukavice}
\textsuperscript{(51)} a \VUgras{cylindr} \textsuperscript{(50)}. Kmotra drží v ruce krásnou
\VUgras{kytici} \textsuperscript{(54)}.

%%K t03-c1-19-1 p
19. --- Za nimi vidíme \VUgras{strýce} \textsuperscript{(56)} a \VUgras{tetu}
\textsuperscript{(58)} dítěte. Teta má elegantní \VUgras{klobouk} \textsuperscript{(59)},
strýc černý \VUgras{redingot} \textsuperscript{(57)} a bílou vestu. Dále jdou
\VUgras{bratranec} \textsuperscript{(60)} a \VUgras{sestřenice} \textsuperscript{(61)}. Ta
drží za ruku \VUgras{sestru} \textsuperscript{(62)} novorozeněte, malou Bertu.

%%K t03-c1-20-1 p
20. --- Bertin druhý \VUgras{bratr} \textsuperscript{(63)} kráčí vzadu. Podává ruku jedné
\VUgras{příbuzné} \textsuperscript{(64)}. \VUgras{Přátelé} \textsuperscript{(65)} rodiny jdou
potom.

%%K t03-tit-7 sub
\VUsaut{4.6mm}
\VUpk{10.2pt}{40}{Diváci.}
\VUsaut{3.4mm}

%%K t03-c1-21-1 p
21. --- Blízko průvodu, zde je \VUgras{žebrák} \textsuperscript{(66)} opřený o svou
\VUgras{berlu} \textsuperscript{(67)}. Prosí o almužnu (žebrá).

%%K t03-c1-22-1 p
22. --- Na druhé straně je malý chlapec velmi \VUgras{zdvořilý} \textsuperscript{(68)}. Zdraví
a přeje «\VUgras{dobrý den}» lidem, které zná.

%%K t03-c1-23-1 p
23. --- Vesničané vyšli ze svých domů, aby viděli křestní průvod. Vidíme
\VUgras{venkovana} \textsuperscript{(69)} s jeho \VUgras{bavlněnou čepicí} a vedle
\VUgras{venkovanku} \textsuperscript{(70)}. Potom \VUgras{stařenu}
\textsuperscript{(71)} opřenou o svou \VUgras{hůl} \textsuperscript{(72)}. Konečně
\VUgras{mlynáře} \textsuperscript{(73)} s jeho velkým \VUgras{plstěným kloboukem}
\textsuperscript{(74)} a mladou venkovanku obutou do \VUgras{dřeváků} \textsuperscript{(75)},
která učí chodit své děťátko.

%%K t03-c1-24-1 p
24. --- Ten výjev je velmi zábavný.

%%K t03-tit-8 sub
\VUsaut{4.0mm}
\VUcentre{12.0pt}{0}{\textit{Druhá scéna.}}
\VUsaut{3.0mm}

%%K t03-tit-9 sub
\VUpk{10.2pt}{40}{Veřejná slavnost.}
\VUsaut{2.4mm}

%%K t03-tit-10 sub
\VUpk{10.2pt}{40}{Vodní hry a závody.}
\VUsaut{3.0mm}

%%K t03-c2-01-1 p
1. --- Nastalo \VUgras{léto}. Žhavé slunce měsíce \VUgras{června} (července nebo
\VUgras{srpna}) pobízí mládež ke koupání a k veslování.

%%K t03-c2-02-1 p
2. --- Na pravém břehu řeky se veslaři svlékají, aby se zúčastnili vodních her a závodů.

%%K t03-c2-03-1 p
3. --- Jeden z nich si zouvá své \VUgras{boty} \textsuperscript{(1)}; rozvazuje je (rozepíná).
Jeho dva kamarádi jsou už hotovi. Od této chvíle mají jen svůj
\VUgras{pletený dres} \textsuperscript{(2)} a \VUgras{plavky} \textsuperscript{(3)}.
Rozmlouvají, zatímco čekají na své přátele. Mluví o pouti a o slavnosti, která se koná na
druhém břehu.

%%K t03-c2-04-1 p
4. --- Na zemi, blízko nich, leží \VUgras{šaty} jejich kamarádů: \VUgras{pár}
\VUgras{polobotek} \textsuperscript{(4)}, \VUgras{boty} \textsuperscript{(5)},
\VUgras{ponožky} \textsuperscript{(6)}, \VUgras{plstěný} \textsuperscript{(7)} a
\VUgras{slaměný} \textsuperscript{(8)} klobouk a \VUgras{vázanka}
\textsuperscript{(9)}.

%%K t03-c2-05-1 p
5. --- Jeden z mladíků pečlivě složil svůj \VUgras{oblek} \textsuperscript{(10)}. Klade jej na
ručník, do něhož jeho soused brzy zabalí oba \VUgras{obleky.}

%%K t03-c2-06-1 p
6. --- Šestý chlapec se opozdil. Má ještě na hlavě svou \VUgras{čepici}
\textsuperscript{(11)}. Nesvlékl si ani svou \VUgras{košili} \textsuperscript{(12)}, ani svůj
\VUgras{límec} \textsuperscript{(13)}, ani své \VUgras{manžety} \textsuperscript{(14)}. Drží v
ruce svou \VUgras{vestu} \textsuperscript{(17)} a má ještě své \VUgras{kalhoty}
\textsuperscript{(16)} a \VUgras{šle} \textsuperscript{(15)}. Jistě nebude hotov na blízký
závod.

%%K t03-c2-07-1 p
7. --- Blízko mladíků vede k břehu řeky pěšinka, po jejíchž stranách je mnoho
\VUgras{bodláků} \textsuperscript{(18)}, kde otec Jakub připravuje mezi \VUgras{rákosím}
\textsuperscript{(19)} \VUgras{ohňostroj} \textsuperscript{(20)}, který se bude odpalovat
večer.

%%K t03-tit-11 sub
\VUsaut{4.4mm}
\VUpk{10.2pt}{40}{Závod.}
\VUsaut{3.4mm}

%%K t03-c2-08-1 p
8. --- Na řece se několik dam, chránících se svými slunečníky, projíždí na
\VUgras{lodce} \textsuperscript{(21)}. Sledují závod, který je velmi zajímavý. V každé
\VUgras{jole} \textsuperscript{(22)} veslují ze všech sil čtyři \VUgras{veslaři}
\textsuperscript{(24)}, zatímco \VUgras{kormidelník} \textsuperscript{(26)} drží
\VUgras{kormidlo} \textsuperscript{(25)} a řídí křehký člun. \VUgras{Vesla}
\textsuperscript{(23)} v taktu bijí do vody a lehké čluny plují závratnou rychlostí.

%%K t03-c2-09-1 p
9. --- Několik mladíků zápasí u pilíře mostu o cenu \VUgras{čelenu} \textsuperscript{(28)}.
Jeden z nich opatrně kráčí po ohebném a kluzkém čelenu, aby dosáhl malé \VUgras{vlaječky}
\textsuperscript{(29)} připevněné na konci. Jeho nešťastný \VUgras{soupeř}
\textsuperscript{(30)} spadl do vody. Plave, aby se dostal na břeh.

%%K t03-c2-10-1 p
10. --- Starší mladíci \VUgras{zápasí na lodích} \textsuperscript{(32)} blízko
\VUgras{nábřeží} \textsuperscript{(33)}. Všem těm hrám přihlíží početné
\VUgras{obecenstvo} \textsuperscript{(31)} opřené o \VUgras{zábradlí}
\VUgras{mostu} \textsuperscript{(27)} a shromážděné na \VUgras{břehu.} Někteří diváci
se však obracejí, aby sledovali hru o \VUgras{stožár s cenou} \textsuperscript{(45)} nebo aby
obdivovali \VUgras{starostův průvod}, který přichází na
\VUgras{rozdílení} cen.

%%K t03-c2-11-1 p
11. --- Nejprve, zde je několik \VUgras{kluků} \textsuperscript{(35)}, kteří jdou před
\VUgras{hudebníky} \textsuperscript{(36)}; potom přichází \VUgras{starosta}
\textsuperscript{(37)}, radní a \VUgras{obecní zastupitelstvo} městečka. Nakonec kráčejí
\VUgras{hasiči} \textsuperscript{(38)}.

%%K t03-tit-12 sub
\VUsaut{5.0mm}
\VUpk{10.2pt}{40}{Pouť.}
\VUsaut{3.6mm}

%%K t03-c2-12-1 p
12. --- Na \VUgras{pouťovém prostranství} \textsuperscript{(39)} je početný zástup lidí. Chodí
se sem dívat na různé \VUgras{boudy}: na \VUgras{dřevěné koně} \textsuperscript{(40)}, na
\VUgras{cirkus} \textsuperscript{(41)}, kde už začalo představení; na \VUgras{veřejný bál}
\textsuperscript{(42)}, kde mladíci a dívky z města požívají rozkoše \VUgras{tance.}

%%K t03-c2-13-1 p
13. --- Vzadu vidíme \VUgras{arénu} \textsuperscript{(44)}, kde statečný španělský
\VUgras{toreador} zápasí se zuřivým \VUgras{býkem} \textsuperscript{(45)}, potom
\VUgras{skluzavku} \textsuperscript{(49)}, \VUgras{střelnici}
\textsuperscript{(46)}, kde se lidé cvičí v zacházení se \VUgras{střelnými zbraněmi} a ve
střelbě do \VUgras{terče.}

%%K t03-c2-14-1 p
14. --- Blízko je \VUgras{pouťové divadlo} \textsuperscript{(47)}, v němž se hraje
\VUgras{veselohra} a \VUgras{drama.} Velmi blízko je bouda \VUgras{obra}
\textsuperscript{(48)} a \VUgras{trpaslíka.} \VUgras{Postava} prvního měří dva metry a patnáct
centimetrů; druhý je sotva tak velký jako dítě.

%%K t03-c2-15-1 p
15. --- Konečně, zde je velký \VUgras{zvěřinec} \textsuperscript{(50)} se svými
\VUgras{divokými zvířaty}, se svým \VUgras{tygrem} \textsuperscript{(51)} s mocnými
\VUgras{drápy}, se svým \VUgras{lvem} \textsuperscript{(52)} s hustou \VUgras{hřívou},
se svými dvěma \VUgras{žirafami} \textsuperscript{(53)} a se svým \VUgras{slonem}
\textsuperscript{(54)}, který tak obratně užívá svého \VUgras{chobotu.}

%%K t03-c2-16-1 p
16. --- \VUgras{Krotitel} \textsuperscript{(55)}, který ta zvířata cvičí, je u dveří boudy.
